\section*{Q1: \texttt{harmonicas.py}}

Nas análises de qualidade de energia em sistemas elétricos trifásicos, é comum modelar distorções nos sinais de tensão e
corrente por meio da decomposição harmônica.
Essa técnica representa sinais periódicos como a soma de componentes senoidais de diferentes frequências, permitindo
identificar e quantificar distúrbios como a distorção harmônica total e a presença de harmônicas específicas que podem
afetar o desempenho de equipamentos e a eficiência do sistema

Além disso, as componentes harmônicas ímpares podem ser classificadas em três categorias, com base na sua ordem $h$
(um número inteiro positivo e ímpar), conforme sua relação com o número de fases do sistema (geralmente três):


\begin{itemize}
\item Sequência zero: quando $h$ é múltiplo de 3, ou seja, $h \mod 3 = 0$.
\item Sequência positiva: quando $h \mod 3 = 1$.
\item Sequência negativa: quando $h \mod 3 = 2$.
\end{itemize}

Implemente uma função que receba uma lista de inteiros positivos representando ordens harmônicas e os
classifique nas três categorias descritas acima.
A função deve retornar uma tupla com três listas:
\begin{itemize}
\item A primeira contendo os valores de sequência zero;
\item A segunda contendo os valores de sequência positiva; e
\item A terceira contendo os valores de sequência negativa.
\end{itemize}


Se não houver valores em alguma categoria, a lista correspondente deve ser vazia.

Se houver valores pares, ignore-os, e não os inclua em nenhuma das três categorias.

Assuma que o argumento é uma lista de inteiros positivos (sem necessitar de uma validação disso interna a sua função).

\begin{minted}{custompython}
def classifica_harmonicas(hs: list) -> tuple[list, list, list]:
    # seu código aqui


# Exemplo de uso
ordem_harmonicas = [5, 7, 9, 3, 2, 11, 13, 15, 1]

zero, pos, neg = classifica_harmonicas(ordem_harmonicas)
print(f"componentes harmônicas de sequência zero: {zero}")
print(f"componentes harmônicas de sequência positiva: {pos}")
print(f"componentes harmônicas de sequência negativa: {neg}")
\end{minted}

Resultado esperado:
\begin{minted}{text}
componentes harmônicas de sequência zero: [9, 3, 15]
componentes harmônicas de sequência positiva: [7, 13, 1]
componentes harmônicas de sequência negativa: [5, 11]
\end{minted}

Dica: Em python, a operação matemática $\mod$ pode ser realizada pelo operador \texttt{\%}.


